\apendice{Especificación de requisitos}

\section{Introducción}

En este anexo se recogen los requisitos principales del editor de diagramas Entidad-Relación desarrollado durante el proyecto. Estos requisitos sirven para describir qué funcionalidades debe ofrecer la aplicación y qué condiciones generales debe cumplir para que su uso sea coherente con los objetivos del trabajo.

El proyecto parte de una aplicación heredada, por lo que los requisitos no se plantean como si el sistema se hubiera diseñado desde cero. Parte del trabajo consistió en analizar la funcionalidad existente, estabilizarla y ampliarla para soportar nuevos elementos del modelo Entidad-Relación. Por este motivo, los requisitos reflejan tanto funcionalidades ya presentes en el editor inicial como ampliaciones incorporadas durante este TFG.

La especificación se centra en los aspectos más relevantes para el alcance del proyecto: edición visual de diagramas, mantenimiento de un modelo interno coherente, validación, transformación al modelo relacional, generación de SQL e importación y exportación de diagramas. No se pretende recoger todos los posibles casos de uso de una herramienta profesional de modelado, sino aquellos que resultan necesarios para el desarrollo realizado.

\section{Objetivos generales}

El objetivo general del sistema es permitir al usuario diseñar diagramas Entidad-Relación mediante una aplicación web y obtener, a partir de ellos, una representación relacional y un script SQL inicial.

A partir de este objetivo general, la aplicación debe permitir crear y modificar elementos del modelo Entidad-Relación, mantener la información del diagrama en una estructura interna coherente, validar el modelo antes de procesarlo y generar una salida que pueda servir como base para la creación de tablas en una base de datos relacional.

También se considera importante que la herramienta permita guardar y recuperar diagramas, de forma que el trabajo realizado por el usuario no dependa únicamente de una sesión concreta del navegador.

\clearpage

\section{Catálogo de requisitos}

En esta sección se presenta un catálogo resumido de los requisitos principales del sistema. Se distinguen requisitos funcionales, relacionados con las operaciones que debe permitir la aplicación, y requisitos no funcionales, relacionados con condiciones generales de uso, mantenimiento y desarrollo.

\begin{table}[H]
    \centering
    \begin{tabular}{p{0.18\textwidth} p{0.70\textwidth}}
        \toprule
        \textbf{Código} & \textbf{Descripción} \\
        \midrule
        RF-01 & El sistema debe permitir crear y editar entidades dentro del lienzo del editor. \\
        RF-02 & El sistema debe permitir añadir atributos a entidades y relaciones. \\
        RF-03 & El sistema debe permitir representar atributos clave, compuestos, multivaluados y discriminantes en los casos soportados por la aplicación. \\
        RF-04 & El sistema debe permitir crear y configurar relaciones entre entidades, incluyendo cardinalidades y roles en los casos soportados por la aplicación. \\
        RF-05 & El sistema debe permitir representar entidades débiles y su dependencia respecto a una entidad propietaria. \\
        RF-06 & El sistema debe mantener un modelo interno coherente con los elementos representados en el lienzo. \\
        RF-07 & El sistema debe validar el diagrama antes de realizar operaciones como la transformación relacional o la generación de SQL. \\
        RF-08 & El sistema debe transformar el diagrama Entidad-Relación validado a una representación relacional. \\
        RF-09 & El sistema debe generar un script SQL a partir de la representación relacional obtenida. \\
        RF-10 & El sistema debe permitir guardar, importar y exportar diagramas en formato JSON. \\
        \bottomrule
    \end{tabular}
    \caption{Catálogo de requisitos funcionales.}
    \label{tab:requisitos-funcionales}
\end{table}

\begin{table}[H]
    \centering
    \begin{tabular}{p{0.18\textwidth} p{0.70\textwidth}}
        \toprule
        \textbf{Código} & \textbf{Descripción} \\
        \midrule
        RNF-01 & La aplicación debe poder ejecutarse en un navegador web moderno. \\
        RNF-02 & La interfaz debe permitir una edición visual clara de los elementos principales del diagrama. \\
        RNF-03 & El código debe mantener una separación razonable entre la interfaz, el modelo interno, la validación y la generación de SQL. \\
        RNF-04 & El sistema debe facilitar la revisión del comportamiento mediante pruebas automáticas y revisión manual. \\
        RNF-05 & La aplicación debe mantener la compatibilidad, en la medida de lo posible, con diagramas guardados en versiones anteriores del modelo interno. \\
        \bottomrule
    \end{tabular}
    \caption{Catálogo de requisitos no funcionales.}
    \label{tab:requisitos-no-funcionales}
\end{table}

% TODO: revisar el catálogo cuando se cierre el soporte final de ISA y agregaciones.

\section{Especificación de requisitos}

Los requisitos funcionales se pueden agrupar en varios bloques. El primero corresponde a la edición visual del diagrama. En este bloque se incluyen las operaciones necesarias para crear entidades, atributos y relaciones, modificar sus nombres y configurar la información asociada a cada elemento.

Un segundo bloque está relacionado con el mantenimiento del modelo interno. Cada elemento representado en el lienzo debe tener una correspondencia en la estructura de datos de la aplicación. Esta estructura es la que permite guardar el diagrama, reconstruirlo posteriormente, validarlo y utilizarlo como entrada para la transformación relacional.

El tercer bloque corresponde a la validación del modelo. Antes de generar una salida relacional o SQL, la aplicación debe comprobar que el diagrama contiene la información mínima necesaria y que no existen configuraciones ambiguas o incompletas. Esta validación no pretende cubrir cualquier posible restricción semántica, pero sí evitar que se generen resultados incoherentes a partir de un diagrama mal definido.

Otro bloque importante es la transformación al modelo relacional. A partir del diagrama validado, el sistema debe obtener una representación basada en tablas, columnas, claves primarias y claves foráneas. Esta representación intermedia permite separar la lógica de transformación de la generación concreta del script SQL.

Finalmente, la aplicación debe permitir generar SQL y trabajar con diagramas guardados. La generación de SQL proporciona una salida textual útil como punto de partida para crear la estructura de una base de datos, mientras que la importación y exportación en JSON permite conservar y recuperar diagramas para continuar su edición.

En cuanto a los requisitos no funcionales, se busca que la aplicación sea utilizable desde un entorno web, que el código sea mantenible y que las distintas partes del sistema estén organizadas de forma que sea posible seguir ampliando el editor sin depender de comportamientos demasiado frágiles.

% TODO: completar o ajustar esta sección si las ampliaciones finales modifican el alcance funcional.